%% ======================================================================
%% docs/update_improve.tex
%% Revised notes for the updated C1--C14 formulation.
%% ======================================================================

\documentclass[11pt,a4paper]{article}

\usepackage{amsmath,amssymb,amsthm}
\usepackage{booktabs}
\usepackage[nameinlink]{cleveref}

\newtheorem{proposition}{Proposition}
\newtheorem{remark}{Remark}

\newcommand{\F}{\mathcal{F}}
\newcommand{\Pset}{\mathcal{P}}
\newcommand{\Aset}{\mathcal{A}}
\newcommand{\Cset}{\mathcal{C}}
\newcommand{\Eset}{\mathcal{E}}
\newcommand{\Mset}{\mathcal{M}}

\title{Event-Based Reformulation of Tail Assignment Constraints C1--C14}
\author{Working Document}
\date{\today}

\begin{document}
\maketitle

\section{Modeling Principle}

The original maintenance model repeats the planning-day index in several
variables even though each flight already has an absolute departure and
arrival timestamp. If all times are measured from the beginning of the
planning horizon, the calendar day of flight $i$ is derived data,
$\lfloor t_i^{\mathrm{dep}}/1440\rfloor$, rather than a decision index.

Maintenance is therefore represented as a virtual flight from an airport to
the same airport. It starts immediately after a scheduled flight, has a fixed
duration, occupies the aircraft, and consumes station capacity. The new binary
variable $z_{ijk}$ selects maintenance type $k$ after flight $i$ for aircraft
$j$. This removes the day dimension from the maintenance variables while
preserving all calendar information in the timestamps.

\section{Update Summary for C1--C14}

The C1--C14 update is based on one modeling decision:
\emph{do not use planning days as decision indices when absolute flight times
already carry the same calendar information.}

The resulting formulation is:

\begin{enumerate}
  \item Keep flight assignment and route sequencing in event time.
  \item Represent each maintenance action as a virtual self-loop leg
  $a\rightarrow a$ with fixed duration $\delta_k$.
  \item Use the maintenance variable
  \[
    z_{ijk}=1
    \iff
    \text{aircraft }j\text{ performs maintenance type }k\text{ after flight }i.
  \]
\end{enumerate}

Thus, maintenance is attached directly to flight events, not to day buckets.
This removes day-indexed binaries while keeping timing, capacity, and
maintenance-feasibility logic fully time-consistent.

\section{A-only ordered-state prototype}

As a compactness experiment, an A-only prototype orders all flights globally
by departure time and replaces the explicit successor variable $w_{i\ell j}$
with one prefix state $q_{rj}$ for each ordered flight position $r$ and
aircraft $j$. This construction is valid only because the A-only state can be
updated while scanning the total time order; route continuity remains enforced
by the corrected basic routing constraints. If flight $i_r$ is not assigned to
$j$, the state is carried forward. If it is assigned without an A check, its
duration is added. If $z_{i_rjA}=1$, the state is reset after that flight.

Let $x_{i_rj}$ be the assignment variable, $z_{i_rjA}$ the A-check event, and
$p_{i_r}$ the flight duration. With $q_{-1,j}=\Phi_{jA}$, a linear
big-$M$ implementation uses:
\begin{align}
q_{rj}&\ge q_{r-1,j}-M_{rj}x_{i_rj},\\
q_{rj}&\le q_{r-1,j}+M_{rj}x_{i_rj},\\
q_{rj}&\ge q_{r-1,j}+p_{i_r}-M_{rj}(1-x_{i_rj})-M_{rj}z_{i_rjA},\\
q_{rj}&\le q_{r-1,j}+p_{i_r}+M_{rj}(1-x_{i_rj})+M_{rj}z_{i_rjA},\\
q_{rj}&\le p_{i_r}z_{i_rjA}+T_A(x_{i_rj}-z_{i_rjA})
  +M_{rj}(1-x_{i_rj}),
\end{align}
where $M_{rj}=\max\{T_A,\Phi_{jA}\}+p_{i_r}$. The state is bounded by
$0\le q_{rj}\le\max\{T_A,\Phi_{jA}\}$, and every assigned flight is
required to satisfy $q_{rj}\le T_A$. In addition, every assigned flight must satisfy
$q_{r-1,j}+p_{i_r}\le T_A$ before a post-flight check can reset the state; a
check after an overdue flight cannot repair the violation. The first two rows carry the state when
the flight is not assigned; the next two rows impose accumulation or relax
the transition at a reset; the final row gives $q_{rj}=p_{i_r}$ at a selected
reset and $q_{rj}\le T_A$ for an assigned non-reset flight.

This ordered-state construction removes $w_{i\ell j}$ only for the A-only
prototype. It must not be generalized to the full hierarchy without adding
separate timestamped C/D event state and proving that route continuity and
maintenance timing remain equivalent. In the current implementation it is
provided as `OrderedAEventMILPScheduler` for controlled parity experiments;
the full validated formulation remains `EventMILPScheduler`.

The same ordered prefix-state construction has been extended to the
hour-based A/B subset as `OrderedABEventMILPScheduler`. It maintains one
$q_{rjA}$ and one $q_{rjB}$ state for each ordered flight position and aircraft.
An A check resets only the A state, while a B check resets both states through
the hierarchy relation. On the clean and near-threshold A/B test instances,
the prototype returns the same optimal objective as the legacy A/B model while
using fewer variables. This is an A/B validation milestone, not yet a claim
for the C/D calendar hierarchy.

The full ordered A/B/C/D prototype has also been tested on the available
feasible ABCD family. On the h=7 instance, the legacy, ordered, and standalone
event models all return the optimal objective 10\,327. The h=15, h=21, and
h=30 files in that family have the same small flight set and all three models
report infeasibility, so they are not treated as horizon-scaling evidence.
The ordered model is therefore ready for broader ABCD testing, but its full
performance claim remains provisional until larger feasible calendar-check
instances are generated.

As a separate fleet-size experiment, one-rotation feasible ABCD instances
were generated for $|\Pc|\in\{5,10,15,20\}$, producing 10, 20, 30, and 40
flights. The reproducible comparison is implemented in
`experiments/compare_ordered_hierarchy.py` and writes
`results/tables/formulation_abcd_fleet_scaling_repro.csv`. Legacy, ordered,
and standalone event models are optimal with identical objectives at all four
fleet sizes. The ordered model uses 47--72\% fewer variables than legacy; its
constraint count is lower for $|\Pc|=5,10$, approximately comparable at 15,
and higher at 20. CPU times are close on this small grid, so these results
support model-size reduction and objective parity, not a universal speedup.

\paragraph{Index choice for maintenance variables.}
The robust event-linked variable is $z_{ijk}$: check type $k$ after flight $i$
for aircraft $j$. A reduced variable such as $z_{jk}$ is not sufficient in the
integrated routing model, because the model must know \emph{after which flight}
the check occurs to enforce turn-time blocking, station capacity over time, and
hour/calendar reset logic. If one wants to remove flight index $i$, an explicit
event-position index (for example, route position $r$) is still required,
which reintroduces sequencing complexity elsewhere.

\section{Sets and Parameters}

Let $\F$ be the scheduled flights, $\Pset$ the aircraft, $\Aset$ the airports,
$\Mset\subseteq\Aset$ the maintenance-capable airports, and
$\Cset=\{A,B,C,D\}$ the check types. Indices $i$ and $\ell$ denote flights,
$j$ denotes an aircraft, and $k$ denotes a check type. Define
$\Cset_H=\{A,B\}$ for flight-hour checks and $\Cset_T=\{C,D\}$ for
calendar-time checks.

For flight $i$, let $o_i$, $a_i$, $t_i^{\mathrm{dep}}$,
$t_i^{\mathrm{arr}}$, and $p_i=t_i^{\mathrm{arr}}-t_i^{\mathrm{dep}}$
be its origin, destination, departure time, arrival time, and flying time.
For aircraft $j$, let $a^0_j$ be its initial airport and $r_j$ its initial
availability time. Let $\Delta$ be the minimum turn time and $\delta_k$ the
fixed duration of check $k$. Maintenance attached to flight $i$ starts at
$\sigma_i=t_i^{\mathrm{arr}}$ and is the virtual leg

\begin{equation}
  m(i,j,k)=\bigl(a_i,a_i,\sigma_i,\sigma_i+\delta_k\bigr).
  \label{eq:virtual_leg}
\end{equation}

Only flights in $\F^M=\{i\in\F:a_i\in\Mset\}$ may be followed by
maintenance. The sparse connection set is

\begin{equation}
  \Eset=\left\{(i,\ell)\in\F\times\F:
  a_i=o_\ell,\quad
  t_\ell^{\mathrm{dep}}\geq t_i^{\mathrm{arr}}+\Delta\right\}.
  \label{eq:connection_set}
\end{equation}

For each check type, define the connections that leave enough time for that
maintenance event:

\begin{equation}
  \Eset_k=\left\{(i,\ell)\in\Eset:
  t_\ell^{\mathrm{dep}}\geq\sigma_i+\delta_k+\Delta\right\}.
  \label{eq:maintenance_connection_set}
\end{equation}

The aircraft-specific set of feasible first flights is

\begin{equation}
  \F^0_j=\left\{i\in\F:o_i=a^0_j,\quad
  t_i^{\mathrm{dep}}\geq r_j\right\}.
  \label{eq:first_flight_set}
\end{equation}

Let $T_k$ be the maximum flying time between qualifying checks for
$k\in\Cset_H$, and let $L_k$ be the maximum elapsed time between qualifying
checks for $k\in\Cset_T$. A limit specified as $d_k$ calendar days is stored as
$L_k=1440d_k$ minutes. Parameters $\Phi_{jk}$ and $\Psi_{jk}$ are the
pre-horizon accumulated flying time and elapsed calendar time, respectively.
The relation $k'\succeq k$ means that check $k'$ is at least as heavy as $k$.

\section{Decision Variables and Objective}

\begin{align}
  x_{ij}&\in\{0,1\},
  &&\text{flight $i$ is assigned to aircraft $j$},
  \label{eq:x_definition}\\
  y_{i\ell j}&\in\{0,1\},
  &&\text{aircraft $j$ operates $\ell$ immediately after $i$},
  \label{eq:y_definition}\\
  z_{ijk}&\in\{0,1\},
  &&\text{aircraft $j$ receives check $k$ immediately after $i$},
  \label{eq:z_definition}\\
  u_{ijk}&\geq0,
  &&\text{flying time accumulated after $i$, before its check},
  \quad k\in\Cset_H.
  \label{eq:u_definition}
\end{align}

Variables $y_{0ij}$ and $y_{i\omega j}$ indicate the first and last flight of
aircraft $j$. The hierarchy-reset indicator is not stored as a variable; it is
the derived expression

\begin{equation}
  q_{ijk}=\sum_{k'\in\Cset:\,k'\succeq k}z_{ijk'}.
  \label{eq:q_definition}
\end{equation}

With assignment costs $c_{ij}$ and maintenance costs $g_{ijk}$, the objective
is

\begin{equation}
  \min
  \sum_{i\in\F}\sum_{j\in\Pset}c_{ij}x_{ij}
  +\sum_{i\in\F^M}\sum_{j\in\Pset}\sum_{k\in\Cset}g_{ijk}z_{ijk}.
  \label{eq:event_objective}
\end{equation}

\section{Complete Event-Based Criteria C1--C14}

\paragraph{C1: Flight coverage.}
Every scheduled flight is assigned to exactly one aircraft:

\begin{equation}
  \sum_{j\in\Pset}x_{ij}=1,
  \qquad \forall i\in\F.
  \tag{C1}\label{eq:event_c1}
\end{equation}

\paragraph{C2: Unique predecessor.}
Every assigned flight has one predecessor or is the first route flight:

\begin{equation}
  \mathbb{1}_{\{\ell\in\F^0_j\}}y_{0\ell j}
  +\sum_{i:(i,\ell)\in\Eset}y_{i\ell j}=x_{\ell j},
  \qquad \forall \ell\in\F,\ j\in\Pset.
  \tag{C2}\label{eq:event_c2}
\end{equation}

Source variables $y_{0\ell j}$ are instantiated only for
$\ell\in\F^0_j$; the indicator in C2 is notation for this sparse domain.

\paragraph{C3: Unique successor.}
Every assigned flight has one successor or is the last route flight:

\begin{equation}
  y_{i\omega j}+\sum_{\ell:(i,\ell)\in\Eset}y_{i\ell j}=x_{ij},
  \qquad \forall i\in\F,\ j\in\Pset.
  \tag{C3}\label{eq:event_c3}
\end{equation}

\paragraph{C4: One time-ordered route per aircraft.}

\begin{align}
  \sum_{i\in\F}y_{0ij}&\leq1,
  &\forall j\in\Pset,\notag\\
  \sum_{i\in\F}y_{i\omega j}&=\sum_{i\in\F}y_{0ij},
  &\forall j\in\Pset.
  \tag{C4}\label{eq:event_c4}
\end{align}

Airport continuity, turn time, and non-overlap are embedded in $\Eset$. Since
all arcs move forward in time, C2--C4 cannot produce a subtour.

\paragraph{C5: Maintenance requires assignment.}

\begin{equation}
  z_{ijk}\leq x_{ij},
  \qquad \forall i\in\F^M,\ j\in\Pset,\ k\in\Cset.
  \tag{C5}\label{eq:event_c5}
\end{equation}

\paragraph{C6: At most one physical check after a flight.}

\begin{equation}
  \sum_{k\in\Cset}z_{ijk}\leq1,
  \qquad \forall i\in\F^M,\ j\in\Pset.
  \tag{C6}\label{eq:event_c6}
\end{equation}

A heavy check satisfies lighter requirements through $q_{ijk}$, so separate
simultaneous A, B, C, and D checks are unnecessary.

\paragraph{C7: Maintenance-airport eligibility.}
Variables $z_{ijk}$ are instantiated only for $i\in\F^M$. A dense equivalent
is

\begin{equation}
  z_{ijk}=0,
  \qquad \forall i\in\F\setminus\F^M,\ j\in\Pset,\ k\in\Cset.
  \tag{C7}\label{eq:event_c7}
\end{equation}

\paragraph{C8: Maintenance duration and the next flight.}
An insufficient connection and a maintenance action cannot both be selected:

\begin{equation}
  y_{i\ell j}+z_{ijk}\leq1,
  \qquad
  \forall(i,\ell)\in\Eset\setminus\Eset_k,\ j\in\Pset,\ k\in\Cset.
  \tag{C8}\label{eq:event_c8}
\end{equation}

For a terminal flight, $z_{ijk}$ is created only if
$\sigma_i+\delta_k$ is within the planning horizon. Thus a multi-day check is
one fixed-duration virtual leg, not one binary variable per occupied day.

\paragraph{C9: Continuous-time station capacity.}
Let $b_a(\theta)$ be the available maintenance positions at airport $a$ and
time $\theta$. Let $\Theta$ contain all possible maintenance start times and
all breakpoints at which $b_a(\theta)$ changes. These are the only times at
which the left-hand side can increase or the right-hand side can decrease:

\begin{equation}
  \sum_{\substack{i\in\F^M:a_i=a\\
                    k\in\Cset:\,\sigma_i\leq\theta<\sigma_i+\delta_k}}
  \sum_{j\in\Pset}z_{ijk}
  \leq b_a(\theta),
  \quad \forall a\in\Mset,\ \theta\in\Theta.
  \tag{C9}\label{eq:event_c9}
\end{equation}

\paragraph{C10: Maintenance hierarchy.}

\begin{equation}
  q_{ijk}=\sum_{k'\succeq k}z_{ijk'},
  \qquad \forall i\in\F^M,\ j\in\Pset,\ k\in\Cset.
  \tag{C10}\label{eq:event_c10}
\end{equation}

C10 is a derived expression and replaces the old daily hierarchy variable.

\paragraph{C11: Flight-hour threshold.}
For A and B requirements, the flight-indexed counter is active only for an
assigned flight and never exceeds the threshold:

\begin{equation}
  p_i x_{ij}\leq u_{ijk}\leq T_kx_{ij},
  \qquad \forall i\in\F,\ j\in\Pset,\ k\in\Cset_H.
  \tag{C11}\label{eq:event_c11}
\end{equation}

\paragraph{C12: Flight-hour accumulation and reset.}
For each route connection, let $R_{\ell k}=T_k+p_\ell$ deactivate propagation
when the successor arc is not selected. The reset term needs only the state
bound $T_k$:

\begin{align}
  u_{\ell jk}
  &\geq u_{ijk}+p_\ell
    -R_{\ell k}(1-y_{i\ell j})-T_kq_{ijk},
    \notag\\
  u_{\ell jk}
  &\leq u_{ijk}+p_\ell
    +R_{\ell k}(1-y_{i\ell j})+T_kq_{ijk},
    \notag\\
  u_{\ell jk}
  &\leq p_\ell+T_k(2-y_{i\ell j}-q_{ijk}),
  \quad
  \forall(i,\ell)\in\Eset,\ j\in\Pset,\ k\in\Cset_H.
  \tag{C12}\label{eq:event_c12}
\end{align}

If $y_{i\ell j}=1$ and no qualifying check follows $i$, the first two
inequalities give $u_{\ell jk}=u_{ijk}+p_\ell$. If a qualifying check follows
$i$, the third inequality and C11 give $u_{\ell jk}=p_\ell$. The counter is
therefore accumulated and reset on the actual aircraft route, with no day-pair
constraints and no endpoint ambiguity.

\paragraph{C13: Pre-horizon accumulated flying time.}
For a first-flight arc, initialize the counter from known aircraft history:

\begin{align}
  u_{ijk}&\geq\Phi_{jk}+p_i-(\Phi_{jk}+p_i)(1-y_{0ij}),
  \notag\\
  u_{ijk}&\leq\Phi_{jk}+p_i+(\Phi_{jk}+p_i)(1-y_{0ij}),
  \quad \forall i\in\F,\ j\in\Pset,\ k\in\Cset_H.
  \tag{C13}\label{eq:event_c13}
\end{align}

C11--C13 replace the old cumulative-flight criterion completely.

\paragraph{C14: Calendar-time spacing without day variables.}
Let $H_0$ and $H_1$ be the horizon boundaries. For $k\in\Cset_T$, the first
deadline is $d^0_{jk}=H_0+L_k-\Psi_{jk}$. If $d^0_{jk}\leq H_1$, require

\begin{equation}
  \sum_{\substack{i\in\F^M:\,H_0\leq\sigma_i\leq d^0_{jk}}}q_{ijk}\geq1,
  \qquad \forall j\in\Pset,\ k\in\Cset_T.
  \tag{C14a}\label{eq:event_c14a}
\end{equation}

Every selected check whose next deadline lies in the horizon must be followed
by another qualifying check within $L_k$ time units:

\begin{equation}
  \sum_{\substack{\ell\in\F^M:\,
                    \sigma_i<\sigma_\ell\leq\sigma_i+L_k}}
  q_{\ell jk}\geq q_{ijk},
  \quad
  \forall i\in\F^M:\sigma_i+L_k\leq H_1,\ j\in\Pset,\ k\in\Cset_T.
  \tag{C14b}\label{eq:event_c14b}
\end{equation}

These timestamp windows require a qualifying check to begin by its due time,
matching the application convention used by the day-indexed model. The virtual
leg then blocks the aircraft for its complete fixed duration. Thus calendar
requirements are preserved without an explicit day set.

\section{Obsolete Variables and Constraints}

The reformulation removes the explicit day index $d$, the old variable
$z_{ijdk}$, the daily maintenance variable $y_{jdk}$, and the daily hierarchy
variable $\gamma_{jdk}$ or $\eta_{jdk}$. It also removes one binary per day of
a multi-day check and all day-pair accumulation constraints.

The symbol $y$ is retained with a new and precise meaning: $y_{i\ell j}$ is a
sparse successor arc. It provides the aircraft sequence required to propagate
the flight-hour counter exactly. The maintenance variable is now only
$z_{ijk}$.

\section{Complexity and Memory Reduction}

Let $n_f=|\F|$, $n_p=|\Pset|$, $n_c=|\Cset|$, $n_d$ be the old number of
days, and $n_e=|\Eset|$. The old maintenance-choice tensor contains
$O(n_fn_pn_dn_c)$ binaries. The new tensor contains only
$O(n_fn_pn_c)$ binaries and is instantiated sparsely over $\F^M$:

\begin{equation}
  O(n_fn_pn_dn_c)\quad\longrightarrow\quad O(n_fn_pn_c).
  \label{eq:maintenance_reduction}
\end{equation}

The complete event-based model contains

\begin{equation}
  O\!\left(n_fn_p+n_en_p+n_fn_pn_c\right)
  \label{eq:event_complexity}
\end{equation}

binary variables, plus $O(n_fn_p|\Cset_H|)$ continuous counters. The sparse
successor variables add $O(n_en_p)$, but $n_e$ is normally much smaller than
$n_f^2$ because airport and time incompatibilities eliminate most arcs. Most
importantly, the maintenance model no longer scales multiplicatively with the
number of planning days.

\begin{table}[ht]
\centering
\caption{Principal variable counts.}
\begin{tabular}{lll}
\toprule
Component & Day-indexed model & Event-based model \\
\midrule
Flight assignment & $n_fn_p$ & $n_fn_p$ \\
Maintenance choice & $n_fn_pn_dn_c$ & $n_fn_pn_c$ \\
Daily maintenance & $n_pn_dn_c$ & removed \\
Daily hierarchy & $n_pn_dn_c$ & derived, not stored \\
Route sequence & implicit & $n_en_p$ sparse binaries \\
Hour state & day-pair constraints & $2n_fn_p$ continuous variables \\
\bottomrule
\end{tabular}
\end{table}

\begin{remark}
The day dimension is removed, but calendar information is not. All timestamps
must use one absolute, horizon-wide time scale; otherwise overnight gaps and
C/D elapsed-time requirements cannot be recovered correctly.
\end{remark}

\section{Model Class: LP, QP, or Convex?}

With binary assignment/connection/maintenance variables
($x_{ij},y_{i\ell j},z_{ijk}\in\{0,1\}$), C1--C14 is a
\textbf{Mixed-Integer Linear Program (MILP)}.

\paragraph{LP.}
The exact integrated model is not a pure LP because of integrality. However,
its LP relaxation (replace $\{0,1\}$ by $[0,1]$) is linear and convex, and is
useful for bounds and diagnostics.

\paragraph{QP.}
A quadratic objective is not required for this reformulation. If quadratic
terms are added (for example, smoothness penalties), the model becomes MIQP,
not simpler than MILP for the core TAP structure.

\paragraph{Convexity.}
The continuous relaxation is convex (polyhedral). The original integer model is
combinatorial and therefore non-convex in decision space.

\paragraph{When can it be pure LP?}
Only under strong simplifications, such as fixed aircraft routes/flight
assignment known in advance, after which maintenance timing may reduce to a
continuous feasibility LP. For integrated assignment plus maintenance timing,
MILP is the correct class.

\section{Implementation Plan for Code and Paper Updates}

This plan updates the codebase and thesis artifacts to the event-based
maintenance reformulation (no day-indexed maintenance variables).

\subsection*{Phase 1: Freeze baseline and define acceptance tests}
\begin{enumerate}
  \item Freeze a baseline commit and keep current regression tests green.
  \item Add acceptance criteria:
  \begin{itemize}
    \item C1--C4 route feasibility unchanged,
    \item maintenance duration blocking via virtual self-loop events,
    \item C/D calendar windows validated by absolute timestamps,
    \item A/B hour reset validated on route replay,
    \item objective decomposition (flight vs maintenance) preserved.
  \end{itemize}
\end{enumerate}

\subsection*{Phase 2: Core model refactor in code}
\begin{enumerate}
  \item Replace day-indexed maintenance tensors by event-linked variables
  $z_{ijk}$ in the new implementation path.
  \item Build sparse connection sets $\Eset$ and $\Eset_k$ from timestamps.
  \item Encode maintenance as fixed-duration virtual self-loop event after
  flight $i$ at airport $a_i$.
  \item Implement C5--C10 in event form (assignment link, single physical
  check, maintenance-station eligibility, duration blocking, hierarchy).
\end{enumerate}

\subsection*{Phase 3: Time-state constraints migration}
\begin{enumerate}
  \item Implement C11--C13 with route-propagated hour-state variable $u$.
  \item Implement C14a/C14b with absolute-time windows for C/D checks.
  \item Remove legacy day-pair C13 logic from the active path after parity
  checks pass.
\end{enumerate}

\subsection*{Phase 4: Validation and parity protocol}
\begin{enumerate}
  \item Unit tests: new tests for maintenance-event placement, C8 blocking,
  C14 deadline enforcement, and hierarchy reset semantics.
  \item Small-instance parity: compare legacy and event models on matched
  feasible instances; verify same optimal objective where semantics overlap.
  \item Independent replay validator: reconstruct route timeline and verify
  all maintenance limits from generated schedules.
  \item Stress runs: scalability checks on h=7/15/21 classes with fixed
  solver/time limits.
\end{enumerate}

\subsection*{Phase 5: Thesis and paper updates}
\begin{enumerate}
  \item Update formal model chapter with C1--C14 event-based equations and
  variable table ($x,y,z,u,q$).
  \item Add a short subsection explaining why $z_{jk}$ is insufficient and
  why event anchor $i$ is required.
  \item Regenerate computational tables/figures and refresh captions to state
  solver settings, instance scope, and parity criteria.
  \item Add reproducibility appendix: exact commands, package versions,
  solver/license notes, and expected outputs.
\end{enumerate}

\subsection*{Phase 6: Packaging and transferability}
\begin{enumerate}
  \item Keep the portable package in \texttt{Thesis/code} as the handoff
  artifact for other devices.
  \item Ensure README runbook includes setup, tests, experiment commands,
  and LaTeX build pipeline.
  \item Add one smoke script per platform (Windows/Linux) for quick
  installation validation.
\end{enumerate}

\paragraph{Exit criteria.}
Migration is complete when: (i) tests pass, (ii) parity/validation reports are
clean, (iii) thesis builds without unresolved references, and (iv) all
reproducibility commands in the package README run on a fresh machine.

\begin{proposition}
Assume that $\Eset$ contains exactly the feasible ordinary connections and C8
is imposed for every maintenance-incompatible connection. Every integral
solution of C1--C14 then defines time-feasible aircraft routes whose selected
maintenance events satisfy duration, station capacity, hierarchy, cumulative
flight-hour limits, and calendar-time limits.
\end{proposition}

\begin{proof}
C1 assigns every flight once. C2--C4 form time-ordered aircraft routes.
C5--C8 attach each maintenance event to an assigned flight at an eligible
station and reserve its duration. C9 enforces station capacity and C10 applies
hierarchy resets. C11--C13 propagate and bound the flight-hour counter along
each selected route. C14 links qualifying calendar-time checks from the known
pre-horizon state through the planning horizon. Hence all stated operating and
maintenance conditions hold.
\end{proof}

\end{document}
